% sec-05: the funding routes opened this week.  Table 25.
\section{The Routes Opened This Week}

The week's letters opened or advanced eleven routes that could pay for part of
the program. Table 25 lists each with its amount, where it stands at the close
of the week, and the one next step it waits on. Four of them pay cash to the
company only after an award the company does not yet hold, and the table says
so rather than leaving a reader to infer it.

\begin{apptable}
\begin{tabularx}{\textwidth}{>{\raggedright\arraybackslash}p{4.4cm} >{\raggedright\arraybackslash}p{3.9cm} >{\raggedright\arraybackslash}p{3.3cm} >{\raggedright\arraybackslash}X}
\headrow \thead{Route} & \thead{Amount} & \thead{At week's close} & \thead{Next step}\\
\midrule
NIH SBIR Phase I & \$306,000; 9 months & Aims to NCI today & A program director \\
NIH SBIR Phase II & \$1,300,000; 24 months & After Phase I & Clinical conduct \\
NSF SBIR Phase I & Up to \$305,000 & Pitch submitted & An invitation \\
NSF SBIR Phase II & Up to \$1,250,000 & After Phase I & None yet \\
NCTI Readiness Pilot & \$1M to \$3.75M each & Not yet eligible & Phase II first \\
Strategic Breakthrough & Up to \$30M, matched & Not yet eligible & Phase II first \\
NAIRR Pilot resources & Compute, not cash & Asked Thursday & A reply \\
Genesis Mission & A subcontributor's place & Offered Wednesday & A reply \\
California Competes & A negotiated tax credit & Asked Monday & Jobs first \\
Contingent payroll earmark & \$40,000 ceiling & Designated & Checked Fridays \\
Private capital & \$5,900,000, in the plan & Ledger at \$0 & The firewall first \\
\end{tabularx}
\end{apptable}
\vspace{-0.60cm}
\tabcap{Table 25. The eleven funding routes the week opened or advanced, each with its amount,\\
where it stands at the close of the week, and the one next step it is waiting on.}

\subsection{The money frame the routes sit inside}

None of the numbers below is re-derived for this packet. The program's direct
cost is \$3,500,000 over five years, or \$700,000 a year, of which \$521,000 is
personnel across 3.95 full-time equivalents \cite{movein2026}. The NIH route of
\$306,000 in Phase I and \$1,300,000 in Phase II totals \$1,606,000, of which
\$1,396,000 is direct, leaving a delta of \$2,104,000 that the route does not
buy; the capitalization plan places \$5,900,000 of private capital behind a
firewall against that delta, a private to federal ratio of 3.67 to 1
\cite{capitalplan2026}. The NSF routes \cite{nsf26510,nsfncti2026} would pay for
the verification work, which the NIH route then cites rather than repeats, so
that no dollar is counted twice. The company's comparable virtual trial work is
\textbf{projected} at \$36,330 per run and about one month, against industry
benchmarks above \$120,000 and 4.5 months \cite{capitalplan2026}; that figure
describes simulation work, not the cost of the clinical trial.

\subsection{The registry every route will ask for}

Every SBIR route in Table 25 will ask for the company's registry record: NIH on
the SBIR/STTR Information form of each application \cite{nihg440}, and NSF before
a full proposal \cite{nsfsbaregistry}. The record is issued by the Small Business
Administration's company registry \cite{sbacompanyreg2026}. Today's form pack
checks its seven fields at the close of a week that brought two applicants and a
private capital question, and records that neither has changed the company's
employee count or ownership, because no offer has been made and no equity sold.
